\chapter{The spin-1 Heisenberg ring}
\label{chap:heisenberg}

This chapter presents our headline application (Lean files
\texttt{CertifiedGaps/Heisenberg/*.lean}): for every even \(L \ge 2304\), the periodic spin-1
antiferromagnetic Heisenberg ring has a unique ground state and
\(\Delta_L > (\log 20)/784 \approx 0.003821\), conditional on two finite seed inequalities. The
seeds are Proposition~5.5 of~\cite{OpenAIHaldane}, reproduced in exact arithmetic in
\texttt{models/spin1/}.

\section{The Hamiltonian}

We define the spin-1 Heisenberg Hamiltonian directly in the Cartesian basis so that the main theorem
statement refers only to standard spin-1 exchange matrices.

\begin{definition}[Spin-1 Heisenberg ring]
  \label{def:heisenberg}
  \lean{CertifiedGaps.SpinOneHeisenberg.leviCivita, CertifiedGaps.SpinOneHeisenberg.spinOne,
    CertifiedGaps.SpinOneHeisenberg.exchange, CertifiedGaps.SpinOneHeisenberg.hamiltonian}
  \leanok
  \uses{def:ring-operators}
  In the Cartesian basis the spin-1 matrices are \((S^\alpha)_{ss'} = -i\,\epsilon_{\alpha s s'}\)
  with the Levi-Civita symbol \(\epsilon\). The exchange is
  \(\mathbf S \cdot \mathbf S = \sum_\alpha S^\alpha \otimes S^\alpha\), and
  \(H_L = \sum_{j} \mathbf S_j \cdot \mathbf S_{j+1}\) on \((\C^3)^{\otimes L}\) with periodic
  boundary conditions.
\end{definition}

\begin{definition}[Spectrum, partition function and purities]
  \label{def:heisenberg-spectrum}
  \lean{CertifiedGaps.SpinOneHeisenberg.hamiltonian_isHermitian,
    CertifiedGaps.SpinOneHeisenberg.spectrum, CertifiedGaps.SpinOneHeisenberg.partitionFunction,
    CertifiedGaps.SpinOneHeisenberg.spatialPurity, CertifiedGaps.SpinOneHeisenberg.thermalPurity}
  \leanok
  \uses{def:heisenberg}
  The Hamiltonian \(H_L\) is Hermitian; \(E_L\) denotes its eigenvalue family,
  \(Z_L(\beta) = \operatorname{Re}\Tr e^{-\beta H_L}\) is its partition function, and \(\Ps\),
  \(\Pt\) are formed from these traces as in Definition~\ref{def:purities}.
\end{definition}

\begin{lemma}
  \label{lem:heisenberg-trace}
  \lean{CertifiedGaps.SpinOneHeisenberg.trace_exp_hamiltonian,
    CertifiedGaps.SpinOneHeisenberg.partitionFunction_eq_sum,
    CertifiedGaps.SpinOneHeisenberg.partitionFunction_spectrum}
  \leanok
  \uses{def:heisenberg-spectrum, lem:ring-trace-exp}
  The matrix trace satisfies \(Z_L(\beta) = \sum_m e^{-\beta E_L(m)}\), matching the partition
  function of the abstract criterion.
\end{lemma}
\begin{proof}
  \leanok
  Apply Lemma~\ref{lem:ring-trace-exp} to the Hermitian matrix \(H_L\).
\end{proof}

\section{The Heisenberg ring is reflection positive}

Because \(iS^\alpha\) is real and antisymmetric in the Cartesian basis, the negated exchange bond is
a sum of squares of real matrices, placing the Heisenberg chain (and its single-ion anisotropic
extension) inside the \texttt{RPRing} class of Chapter~\ref{chap:rings}.

\begin{definition}[Real spin-1 generators]
  \label{def:spin-one-generators}
  \lean{CertifiedGaps.SpinOne.G, CertifiedGaps.SpinOne.G_transpose}
  \leanok
  \uses{def:heisenberg}
  Each generator \(G_\alpha = i S^\alpha\) has real antisymmetric entries \((G_\alpha)_{ss'} =
  \epsilon_{\alpha s s'}\), and \(-\mathbf S \cdot \mathbf S = \sum_\alpha G_\alpha \otimes
  G_\alpha\).
\end{definition}

\begin{definition}[The anisotropic spin-1 ring]
  \label{def:spin-one-large-d}
  \lean{CertifiedGaps.SpinOneLargeD.Sz2, CertifiedGaps.SpinOneLargeD.model}
  \leanok
  \uses{def:rpring, def:spin-one-generators}
  For \(D \in \R\), the anisotropic spin-1 ring is the reflection-positive ring with bond factors
  \(C_\alpha = G_\alpha\), on-site term \(o = D\,(S^z)^2 = D \operatorname{diag}(1, 1, 0)\), and
  internal reflection \(F = 1\).
\end{definition}

\begin{lemma}
  \label{lem:large-d-heisenberg}
  \lean{CertifiedGaps.SpinOneLargeD.bondMatrix_eq, CertifiedGaps.SpinOneLargeD.oc_eq,
    CertifiedGaps.SpinOneLargeD.hamiltonian_eq, CertifiedGaps.SpinOneLargeD.partitionFunction_zero}
  \leanok
  \uses{def:spin-one-large-d, def:heisenberg-spectrum}
  The bond of this ring is \(\mathbf S \cdot \mathbf S\) and its Hamiltonian is
  \(\sum_j \mathbf S_j \cdot \mathbf S_{j+1} + D \sum_j (S^z_j)^2\). At \(D = 0\) its partition
  functions are those of the Heisenberg ring.
\end{lemma}
\begin{proof}
  \leanok
  The bond identity \(-\sum_\alpha G_\alpha \otimes G_\alpha = -\sum_\alpha (iS^\alpha) \otimes
  (iS^\alpha) = \sum_\alpha S^\alpha \otimes S^\alpha\) holds entrywise from \(i^2 = -1\). At
  \(D = 0\) the on-site terms vanish.
\end{proof}

\begin{theorem}[Spatial moment inequalities for the Heisenberg ring]
  \label{thm:heisenberg-moments}
  \lean{CertifiedGaps.SpinOneHeisenberg.spatialMomentBounds}
  \leanok
  \uses{lem:large-d-heisenberg, thm:rp-moment-bounds}
  For every \(\beta > 0\), \(L \mapsto Z_L(\beta)\) satisfies the spatial moment inequalities
  above \(2\).
\end{theorem}
\begin{proof}
  \leanok
  Apply Theorem~\ref{thm:rp-moment-bounds} to the ring of Definition~\ref{def:spin-one-large-d} at
  \(D = 0\), whose partition functions are the Heisenberg ones by
  Lemma~\ref{lem:large-d-heisenberg}.
\end{proof}

\section{The gap from thermal-MPO inputs}

We package the spatial moment inequalities and the two seed purity bounds into a single hypothesis
record \texttt{ThermalMPOInputs} and specialize the symmetric doubling theorem
(Theorem~\ref{thm:gap-of-seeds}) to obtain the gap bound \(\log(1/(5\varepsilon_0))/\beta_0\).

\begin{definition}[Thermal-MPO inputs]
  \label{def:thermal-mpo-inputs}
  \lean{CertifiedGaps.SpinOneHeisenberg.ThermalMPOInputs}
  \leanok
  \uses{def:heisenberg-spectrum, def:moment-bounds}
  The structure \(\mathrm{Inputs}(n_0, \beta_0, \varepsilon_0)\) packages the spatial moment
  inequalities above \(n_0\) at every \(\beta > 0\) together with the seeds
  \(\Ps(n_0, \beta_0) \ge 1 - \varepsilon_0\) and \(\Pt(2n_0, \beta_0) \ge 1 - \varepsilon_0\).
\end{definition}

\begin{theorem}[Gap from inputs]
  \label{thm:heisenberg-gap-of-inputs}
  \lean{CertifiedGaps.SpinOneHeisenberg.gapGT_of_inputs,
    CertifiedGaps.SpinOneHeisenberg.gapGT_log_twenty}
  \leanok
  \uses{def:thermal-mpo-inputs, thm:gap-of-seeds, lem:heisenberg-trace}
  Let \(n_0 > 0\) be even, \(\beta_0 > 0\) and \(0 < \varepsilon_0 \le 1/100\). Under
  \(\mathrm{Inputs}(n_0, \beta_0, \varepsilon_0)\), every even \(L \ge n_0\) satisfies
  \(\GapGT\bigl(E_L, \log(1/(5\varepsilon_0))/\beta_0\bigr)\). In particular
  \(\mathrm{Inputs}(2304, 784, 1/100)\) gives \(\GapGT(E_L, \log 20 / 784)\) for even
  \(L \ge 2304\).
\end{theorem}
\begin{proof}
  \leanok
  Apply Theorem~\ref{thm:gap-of-seeds}, identifying the partition functions by
  Lemma~\ref{lem:heisenberg-trace}.
\end{proof}

\section{The seeds and the main theorem}

Combining the unconditional spatial moment inequalities (Theorem~\ref{thm:heisenberg-moments}) with
the two seed inequalities at \((n_0, \beta_0) = (2304, 784)\) yields the uniform Haldane gap bound.

\begin{lemma}[Spatial seed]
  \label{lem:heisenberg-seed-spatial}
  \lean{CertifiedGaps.SpinOneHeisenberg.seed_spatial}
  \leanok
  \uses{def:heisenberg-spectrum}
  The spatial purity at \((2304, 784)\) satisfies \(\Ps(2304, 784) \ge 99/100\).
\end{lemma}
\begin{proof}
  This is Proposition~5.5 of~\cite{OpenAIHaldane}, reproduced in exact integer and rational
  arithmetic in \texttt{models/spin1/}; see also~\cite{CertificateNote}. It is not yet proved in
  Lean, as it is a finite numerical statement about rings of \(2304\) and \(4608\) sites.
\end{proof}

\begin{lemma}[Thermal seed]
  \label{lem:heisenberg-seed-thermal}
  \lean{CertifiedGaps.SpinOneHeisenberg.seed_thermal}
  \leanok
  \uses{def:heisenberg-spectrum}
  The thermal purity at \((4608, 784)\) satisfies \(\Pt(4608, 784) \ge 99/100\).
\end{lemma}
\begin{proof}
  This is certified alongside Lemma~\ref{lem:heisenberg-seed-spatial} and is not yet proved in Lean.
\end{proof}

\begin{theorem}[The spin-1 Heisenberg gap]
  \label{thm:haldane-gap}
  \lean{CertifiedGaps.SpinOneHeisenberg.thermalMPOInputs,
    CertifiedGaps.SpinOneHeisenberg.spin_one_heisenberg_gap}
  \leanok
  \uses{thm:heisenberg-moments, lem:heisenberg-seed-spatial, lem:heisenberg-seed-thermal,
    thm:heisenberg-gap-of-inputs}
  For every even \(L \ge 2304\), the periodic spin-1 Heisenberg ring has a unique ground state,
  and every other eigenvalue lies more than \((\log 20)/784 \approx 0.003821\) above the ground
  energy.
\end{theorem}
\begin{proof}
  \leanok
  Theorem~\ref{thm:heisenberg-moments} and the two seeds give
  \(\mathrm{Inputs}(2304, 784, 1/100)\); apply Theorem~\ref{thm:heisenberg-gap-of-inputs}.
\end{proof}

\begin{remark}
  In Lean, \texttt{\#print axioms spin\_one\_heisenberg\_gap} reports \texttt{sorryAx} only through
  the two seed lemmas. Smaller seed routes for a kernel-checked proof of the spin-1 seed are studied
  in \texttt{models/lean\_seeds/}; see Chapter~\ref{chap:seedkernel}.
\end{remark}
